\section{Introduction}
\label{sec:introduction}

Ransomware compresses the defender's decision window. A detector must identify a
malicious program early enough to prevent encryption and lateral impact, yet much
of the most discriminative evidence appears only after execution. Static analysis
can score a Windows portable executable (PE) before it runs and is therefore
attractive for low-latency screening. Modern static detectors learn from byte
distributions, header fields, imports, sections, strings, and related PE
attributes; tree ensembles remain strong tabular baselines on EMBER-style
representations \cite{anderson2018ember,harang2020sorel,ke2017lightgbm}. Their
speed, however, does not remove brittleness under packing, obfuscation, family or
temporal shift, and manipulation of non-functional PE content
\cite{ucci2019survey,gibert2020rise,arp2022dos}.

Dynamic analysis observes execution through file-system, process, registry,
memory, and network events. It can reveal destructive intent that a static
representation conceals and is especially relevant to ransomware, whose impact
is expressed through sequences of discovery, access, modification, and
encryption \cite{alrimy2018survey,kharraz2015cutting,bacci2021ransomware}.
Dynamic evidence nevertheless has its own failure modes. It costs time and
compute, its coverage depends on the analysis environment, and environment-aware
or sleep-delayed malware may withhold behavior during the observation window
\cite{sgandurra2016automated,sikorski2012practical}. Hybrid detection is thus a
natural design choice, but an unqualified average assumes that the reliability of
static and dynamic scores is fixed. This is hard to justify when behavioral
evidence arrives sequentially and may change in informativeness across a session.

The static model also creates two operational problems. First, continuous
adaptation requires labels for emerging, obfuscated, or ambiguous samples.
Expert attention is limited, so the system should prioritize examples that are
informative for retraining rather than annotate an arbitrary stream. Classical
uncertainty sampling selects points close to the current decision boundary
\cite{settles2009active,cohn1996active}, but output ambiguity alone does not reveal
whether the model relies on internally conflicting or concentrated evidence.
Second, an explanation graphic is not proof that the identified features govern
the prediction. Attribution methods can be unstable, can be evaluated using
unrealistic perturbations, and can encourage causal interpretations unsupported
by a predictive model \cite{guidotti2018survey,yeh2019infidelity,
warnecke2020evaluating}. In malware analysis this concern is acute because PE
features are constrained and correlated: changing an isolated coordinate in a
vector need not correspond to any executable file.

This paper studies the intersection of those problems through a multi-stage
framework centered on a LightGBM static classifier and native TreeSHAP
attributions. The framework includes an uncertainty gate for allocating dynamic
analysis, a stateful PID rule for fusing aligned static and dynamic evidence, an
active-learning acquisition score derived from prediction and explanation
signals, and a controlled fidelity protocol. The work deliberately separates
three claims that are often blended: predictive performance, label-selection
quality, and explanation fidelity. Each requires different controls. Moreover,
the available artifact supports quantitative conclusions for active learning and
feature-space fidelity, but not for PID fusion, because it does not contain
independently observed dynamic scores aligned with the static samples. Treating
that absence as a result would turn a system design into an empirical claim.

We organize the investigation around three research questions:

\begin{itemize}
    \item \textbf{RQ1:} At a fixed training-label budget, does a query score that
    combines predictive uncertainty with signed TreeSHAP geometry improve ranking
    performance over random and uncertainty-only acquisition?
    \item \textbf{RQ2:} Under an explicitly defined feature-space intervention,
    does deleting the largest absolute TreeSHAP contributions change model
    behavior more than random deletion with the same coordinate budget?
    \item \textbf{RQ3:} What evidence and controls are required to determine
    whether stateful PID fusion improves over static-only, dynamic-only, fixed,
    and simpler P/PI rules without confusing fusion time with dynamic-analysis or
    end-to-end latency?
\end{itemize}

The principal contributions are as follows.

\begin{itemize}
    \item \textbf{A provenance-aware multi-stage specification.} We formulate a
    static--dynamic architecture and a discrete PID update whose state is reset at
    independent session boundaries. The evaluation protocol locks gains,
    thresholds, and fusion baselines on validation data, requires one-to-one
    evidence linkage, and distinguishes offline score replay from signal
    acquisition and end-to-end latency. This contribution is a falsifiable design;
    no PID performance value is reported without aligned real observations.

    \item \textbf{An explanation-guided acquisition policy with controlled
    ablations.} The query score combines normalized predictive uncertainty,
    directional entropy of positive versus negative TreeSHAP mass, and variance
    across signed feature contributions. Under identical initial sets, splits,
    budgets, and five paired seeds, the full policy reached
    $0.9749\pm0.0040$ ROC--AUC with 700 training labels. The AUC difference from
    uncertainty sampling was positive in all five seeds. We report AP, FPR, TPR,
    partial ablations, learning curves, and query-time scope rather than reducing
    the outcome to a single favorable metric.

    \item \textbf{A controlled fidelity assessment for PE feature vectors.}
    Native TreeSHAP additivity is checked on raw log-odds, selected coordinates
    are replaced by arithmetic means estimated only from the training partition,
    and top-$|\mathrm{SHAP}|$ deletion is compared with 20 uniform random masks at
    each budget. At $k=25$, the top-SHAP AUC drop was 0.1616, approximately 630
    times the mean random drop of 0.000257. Faithfulness, explanation consistency,
    and contrastive separation are reported separately because they have
    different scales and assumptions; an Overall Fidelity Score (OFS) is not
    constructed without a validated aggregation rule.

    \item \textbf{A replay-oriented evidence package.} The measured run records
    fixed splits, source hashes, 150 model checkpoints, query decisions,
    per-sample predictions, mask indices, and pre/post-mask scores. Independent
    verification replayed 120,000 active-learning predictions, 12,500 query
    records, 63 mask batches, and 25,200 fidelity predictions. This supports
    computational traceability while explicitly not claiming provenance that the
    source CSV does not provide.
\end{itemize}

The evaluation is intentionally described as a row-disjoint internal pilot. The
CSV contains 2,000 rows and 2,381 numeric coordinates, but the repository does
not establish the original PE identities, collection dates, malware families, or
whether every positive label specifically denotes ransomware. Consequently, the
paper does not call the experiment a SOREL-20M benchmark, a family-held-out
test, a temporal test, or a zero-day evaluation. These limits reduce the scope of
the conclusion, but they also make the evidence auditable and prevent an
attractive system narrative from outrunning its data.

The remainder of this paper is structured as follows. Section~\ref{sec:related}
positions the work relative to static/dynamic detection, active learning, and
security XAI. Section~\ref{sec:problem} defines the problem, threat model, and
claim boundaries. Section~\ref{sec:method} presents the architecture, PID
specification, acquisition policy, and fidelity measures. Sections~\ref{sec:exp}
and~\ref{sec:results} describe the reproducible protocol and measured outcomes.
Sections~\ref{sec:discussion} and~\ref{sec:threats} discuss implications and
validity threats, and Section~\ref{sec:conclusion} concludes the paper.
